% ---------------------------------------------------------------- front matter
\begin{center}
{\sffamily\small RESEARCH ARTICLE}\\[6pt]
{\sffamily\bfseries\LARGE Doubly Nested Semantic Communication: Rate-Nested\\[2pt] Token Transmission and Time-Nested Language-Model\\[2pt] Memories for Drift-Resilient 6G Links\par}
\vspace{10pt}
{\normalsize Abraheem Rashid$^{1}$ \,|\, Faisal Iradat$^{2}$ \,|\, Waseem Iqbal$^{3}$ \,|\, Yawar Abbas Bangash$^{4}$\par}
\vspace{6pt}
{\footnotesize
$^{1}$Department of Computer Science, Brunel University London, London, United Kingdom\\
$^{2}$Department of Computer Science, School of Mathematics and Science, Institute of Business Administration, Karachi, Pakistan\\
$^{3}$Department of Electrical and Computer Engineering, Sultan Qaboos University, Muscat, Oman\\
$^{4}$School of Computing and Artificial Intelligence, FET, Sunway University, Selangor Darul Ehsan, Malaysia\par}
\vspace{6pt}
{\footnotesize
\textbf{Correspondence:} Waseem Iqbal, Department of Electrical and Computer Engineering, Sultan Qaboos University, Muscat, Oman. Email: \href{mailto:m.waseem@squ.edu.om}{m.waseem@squ.edu.om}\\[2pt]
\textbf{Email addresses:} 2546393@brunel.ac.uk (A. Rashid); firadat@iba.edu.pk (F. Iradat); m.waseem@squ.edu.om (W. Iqbal); yawarb@sunway.edu.my (Y. A. Bangash)\par}
\end{center}

\vspace{6pt}
\noindent\rule{\textwidth}{0.4pt}

\noindent\textbf{\sffamily Abstract}\par\smallskip
\noindent Large language models (LLMs) enable a semantic receiver to recover text from few channel symbols, because the model shared by both ends predicts what the transmitter is likely to send. After deployment, however, the content drifts, and the air time needed to keep the receiver's knowledge aligned is rarely accounted for. This article proposes doubly nested semantic communication (DN-SC), a token-level transceiver that is nested both in rate and in time. Each token is transmitted as a prefix of a 16-symbol codeword whose shells are power-controlled separately, so that the receiver can request further symbols only for uncertain tokens; a frozen GPT-2 model serves as the prior of a Bayesian beam-search decoder; and three receiver memories, namely the conversation context, a logit bias and low-rank adapter slots, adapt the prior at three timescales. The return-link requests and every synchronized bit are converted into channel symbols. On a hierarchical drift stream with six seeds (AWGN, mean signal-to-noise ratio of 10\,dB), DN-SC attains 90\,\% token accuracy with 7.08 real channel symbols per token, compared with 9.23 for a frozen fixed-rate receiver with the same decoder. This corresponds to a reduction of 23.3\,\% with the requests included, observed on every seed, and of 27.9\,\% (95\,\% confidence interval 21.6--31.7\,\%) under free-running decoding. The context memory, which requires no synchronization, provides about three quarters of the gain of the memories, whereas the shipped memories recover their synchronization cost relative to it only when one update serves more than approximately 6100 messages.

\medskip
\noindent\textbf{\sffamily Keywords}\par\smallskip
\noindent 6G \,|\, semantic communication \,|\, large language models \,|\, joint source--channel coding \,|\, incremental redundancy \,|\, continual adaptation \,|\, nested learning

\noindent\rule{\textwidth}{0.4pt}
\vspace{4pt}
